\chapter{Axillary Node Clearance}

\section*{What Is This Surgery?}
Axillary node clearance, also known as axillary lymph node dissection (ALND), is a surgical procedure involving the systematic removal of lymph nodes and associated fibrofatty tissue from the armpit, or axilla. This operation is primarily performed in the context of breast cancer management when there is confirmed or highly suspected spread of cancer cells to the regional lymph nodes. The lymphatic system serves as a common pathway for cancer dissemination, and the axillary nodes represent the primary drainage basin for breast tissue. By excising these nodes, surgeons aim to achieve local disease control, accurately stage the cancer, and provide crucial prognostic information that guides subsequent adjuvant therapies such as chemotherapy, radiation, and endocrine treatment. While once a routine component of breast cancer surgery, its application has become more selective with advancements in sentinel lymph node biopsy techniques.

\section*{Alternative Names}
\begin{itemize}
  \item Axillary lymph node dissection (ALND)
  \item Axillary dissection
  \item Axillary lymphadenectomy
  \item Complete axillary dissection
  \item Level I and Level II axillary lymphadenectomy
  \item Lymph node clearance of the armpit
\end{itemize}

\section*{Relevant Surgical Anatomy}
The axilla is a pyramidal space located between the upper arm and the chest wall, serving as a critical conduit for neurovascular structures and a reservoir for lymph nodes draining the breast and upper limb. Its boundaries are surgically important: anteriorly by the pectoralis major and minor muscles, posteriorly by the latissimus dorsi, teres major, and subscapularis muscles, medially by the serratus anterior muscle and the upper ribs, and laterally by the humerus. The apex of the axilla points towards the cervicoaxillary canal, while the base is formed by the axillary fascia and skin.

Within this fibrofatty space lie the axillary artery and vein, the cords and branches of the brachial plexus, and typically 10-50 axillary lymph nodes. These nodes are surgically classified into three levels relative to the pectoralis minor muscle:
\begin{itemize}
  \item \textbf{Level I:} Nodes lateral to the lateral border of the pectoralis minor.
  \item \textbf{Level II:} Nodes deep to (or posterior to) the pectoralis minor.
  \item \textbf{Level III:} Nodes medial to the medial border of the pectoralis minor.
\end{itemize}
Axillary node clearance typically involves removal of Level I and Level II nodes. Key nerves that must be identified and meticulously preserved during dissection include the long thoracic nerve (supplying the serratus anterior muscle, injury leads to winged scapula), the thoracodorsal nerve (supplying the latissimus dorsi muscle), and the intercostobrachial nerve (a sensory nerve providing sensation to the medial arm and axilla, often sacrificed or injured leading to numbness). The axillary vein forms the superior boundary of the dissection field, and its branches, such as the lateral thoracic vein, must be carefully ligated.

\section*{Surgical Principle \& Rationale}
The fundamental principle of axillary node clearance is based on the understanding of breast cancer's metastatic pathway, where lymphatic spread to the regional axillary lymph nodes is a common and prognostically significant event. The rationale for this surgery is twofold: therapeutic and staging. Therapeutically, it aims to remove all macroscopic and microscopic cancer cells within the axillary lymph nodes, thereby achieving local disease control and reducing the risk of regional recurrence. From a staging perspective, the procedure provides a definitive assessment of the nodal burden, including the number of involved nodes, the presence of extracapsular extension, and the size of metastatic deposits.

The operative concept involves the meticulous dissection and removal of the fibrofatty tissue containing the Level I and II axillary lymph nodes, while carefully preserving vital neurovascular structures. The technical goals are to achieve a complete clearance of all visible nodal disease, obtain an accurate pathological specimen for precise staging, and minimize morbidity by preventing injury to the long thoracic, thoracodorsal, and intercostobrachial nerves, as well as the axillary vessels. This comprehensive approach ensures optimal local control and provides critical information for tailoring adjuvant systemic and radiation therapies, ultimately impacting patient prognosis and survival.

\section*{History \& Surgical Evolution}
The management of the axilla in breast cancer has undergone significant evolution over the past century. Early pioneers like William Halsted, in the late 19th and early 20th centuries, advocated for radical mastectomy, which included an en bloc resection of the breast, pectoralis muscles, and extensive axillary lymph node dissection (including Level III nodes). This aggressive approach was based on the "centrifugal spread" theory of cancer and aimed for maximal local control.

Throughout the mid-20th century, modified radical mastectomy, which preserved the pectoralis muscles, became the standard, still incorporating a comprehensive axillary dissection. A pivotal shift occurred in the 1990s with the introduction and widespread adoption of sentinel lymph node biopsy (SLNB) by surgeons such as Donald Morton and Armando Giuliano. SLNB allowed for the identification and removal of only the first draining lymph nodes, sparing patients with negative sentinel nodes from the morbidity of a full axillary clearance. This innovation dramatically reduced the incidence of lymphedema and other complications. Further refinement came with the ACOSOG Z0011 trial results in 2011, which demonstrated that in specific patients with T1/T2 breast cancer, 1-2 positive sentinel nodes, and planned whole-breast radiation, a complete axillary dissection did not improve survival outcomes compared to SLNB alone. This landmark study further narrowed the indications for axillary node clearance, reserving it for patients with a higher burden of nodal disease or specific clinical scenarios.

\section*{Clinical Indications}
Axillary node clearance is indicated in specific clinical scenarios where the axillary lymph nodes are known or highly suspected to be involved with breast cancer. The primary indications include:

\begin{itemize}
  \item \textbf{Clinically positive axillary lymph nodes:} Detected by physical examination, imaging (ultrasound, MRI, PET/CT), or fine-needle aspiration/core biopsy confirming malignancy, especially if not downstaged by neoadjuvant chemotherapy.
  \item \textbf{Positive sentinel lymph node biopsy (SLNB) with specific criteria:} When SLNB reveals multiple positive nodes (typically three or more), or when there is extracapsular extension, particularly in patients not receiving whole-breast radiation or those undergoing mastectomy.
  \item \textbf{Residual nodal disease after neoadjuvant systemic therapy:} For patients who initially presented with clinically positive axillary nodes and received chemotherapy before surgery, but still have evidence of residual cancer in the axilla after treatment.
  \item \textbf{Inflammatory breast cancer:} A highly aggressive form of breast cancer where nodal involvement is common and comprehensive axillary staging is crucial for treatment planning.
  \item \textbf{Locally advanced breast cancer:} Cases where the tumor is large or has spread extensively to regional lymph nodes, necessitating aggressive local control.
  \item \textbf{Recurrent disease in the axilla:} In selected cases of isolated axillary recurrence after previous breast cancer treatment, where re-excision is feasible and beneficial.
  \item \textbf{Inability to perform sentinel lymph node biopsy:} Rarely, if SLNB fails or is technically not feasible, and there is a high suspicion of nodal involvement.
\end{itemize}

\section*{Clinical Positioning}
Axillary node clearance is primarily positioned as a secondary or definitive treatment step in the multidisciplinary management of breast cancer, rather than a first-line diagnostic or therapeutic procedure on its own. Its role has become highly selective following the widespread adoption of sentinel lymph node biopsy (SLNB).

\begin{itemize}
  \item \textbf{Secondary/Adjunct Role:} In the majority of early-stage breast cancer patients, SLNB is the initial axillary staging procedure. If the sentinel nodes are negative for cancer, a full axillary clearance is typically avoided to prevent associated morbidity. Thus, ALND serves as a secondary procedure, performed only if SLNB confirms nodal involvement beyond specific criteria (e.g., more than 1-2 positive nodes, or in patients not meeting ACOSOG Z0011 criteria).
  \item \textbf{Primary/Definitive Treatment Role:} For patients who present with clinically obvious or biopsy-proven extensive axillary nodal disease (e.g., multiple palpable nodes, large nodal masses, or extensive disease after neoadjuvant therapy), ALND may be considered a primary therapeutic intervention to achieve local disease control. In these cases, it is often performed concurrently with mastectomy or lumpectomy.
  \item \textbf{Salvage Role:} In rare instances of isolated axillary recurrence after prior treatment, ALND may be performed as a salvage procedure to remove recurrent disease.
\end{itemize}
The decision to proceed with axillary node clearance is made after careful consideration of the patient's overall cancer stage, tumor biology, response to neoadjuvant therapy, and potential for morbidity, always within a multidisciplinary team framework.

\section*{Surgical Technique \& Procedure}
Axillary node clearance is performed under general anesthesia, typically in conjunction with a lumpectomy or mastectomy.

\begin{itemize}
  \item \textbf{Anesthesia and Positioning:} The patient is placed in a supine position on the operating table. The ipsilateral arm is abducted to approximately 90 degrees and externally rotated, supported on an arm board, to provide optimal exposure of the axilla. The chest, shoulder, and arm are prepped and draped in a sterile fashion.
  \item \textbf{Incision:} A transverse or oblique incision is made in the axillary crease, often 5-8 cm in length, to ensure good cosmetic outcome and adequate exposure. Alternatively, the incision may be an extension of a mastectomy incision. Skin flaps are raised superiorly and inferiorly to expose the axillary contents.
  \item \textbf{Identification of Landmarks:} The lateral border of the pectoralis major muscle is identified. The axillary vein, which forms the superior boundary of the dissection, is carefully identified and skeletonized.
  \item \textbf{Dissection of Axillary Contents:} The fibrofatty tissue containing the lymph nodes is meticulously dissected from the chest wall (medially), the latissimus dorsi and subscapularis muscles (posteriorly), and the brachial plexus (superiorly). The dissection proceeds from lateral to medial, sweeping the nodal tissue towards the chest wall.
  \item \textbf{Nerve Preservation:} Critical nerves are identified and preserved:
    \begin{itemize}
      \item The \textbf{long thoracic nerve} (Nerve of Bell) is identified running vertically on the surface of the serratus anterior muscle along the medial chest wall.
      \item The \textbf{thoracodorsal nerve and vessels} are identified running obliquely on the surface of the latissimus dorsi muscle, often posterior to the axillary vein.
      \item The \textbf{intercostobrachial nerve}, a sensory nerve, is often encountered crossing the axilla. While efforts are made to preserve it, it may be sacrificed if it obstructs complete nodal clearance or is intimately involved with tumor.
    \end{itemize}
  \item \textbf{Node Removal (Levels I and II):} The entire fibrofatty contents of Level I (lateral to pectoralis minor) and Level II (deep to pectoralis minor) are removed en bloc. The pectoralis minor muscle itself is usually preserved. The specimen is carefully oriented and sent for pathological examination.
  \item \textbf{Hemostasis and Irrigation:} Meticulous hemostasis is achieved using electrocautery and ligatures to prevent hematoma formation. The wound is irrigated with saline.
  \item \textbf{Drain Placement:} A closed-suction drain (e.g., Jackson-Pratt or Blake drain) is typically placed in the axillary dead space to prevent seroma formation. The drain is brought out through a separate stab incision and secured to the skin.
  \item \textbf{Closure:} The skin flaps are reapproximated, and the wound is closed in layers using absorbable sutures for the subcutaneous tissue and non-absorbable sutures or staples for the skin. A sterile dressing is applied.
\end{itemize}

\section*{Preoperative Workup \& Preparation}
Thorough preoperative workup and preparation are essential to ensure patient safety and optimize surgical outcomes for axillary node clearance.

\begin{itemize}
  \item \textbf{Comprehensive Medical Evaluation:} A complete history and physical examination are performed to assess the patient's overall health, identify comorbidities, and evaluate fitness for general anesthesia. This includes cardiac and pulmonary assessments, especially in older patients or those with pre-existing conditions.
  \item \textbf{Laboratory Investigations:} Routine blood tests typically include a complete blood count (CBC), basic metabolic panel (BMP), liver function tests (LFTs), and coagulation studies (PT/INR, PTT). Blood typing and cross-matching may be performed if significant blood loss is anticipated, though it is rare for ALND alone.
  \item \textbf{Imaging Studies:} Preoperative imaging, such as mammography, breast ultrasound, and often breast MRI, is crucial to define the extent of the primary breast tumor. Axillary ultrasound with targeted biopsy of suspicious nodes is standard to confirm nodal involvement. PET/CT scans may be used for systemic staging in higher-risk patients.
  \item \textbf{Medication Review and Adjustment:} All medications, particularly anticoagulants (e.g., warfarin, DOACs) and antiplatelet agents (e.g., aspirin, clopidogrel), are reviewed. These may need to be held or bridged according to institutional protocols to minimize bleeding risk.
  \item \textbf{NPO Status:} Patients are instructed to be NPO (nil per os) for at least 6-8 hours prior to surgery to reduce the risk of aspiration during anesthesia.
  \item \textbf{Prophylactic Antibiotics:} Intravenous prophylactic antibiotics (e.g., cefazolin) are administered within 60 minutes prior to incision to reduce the risk of surgical site infection.
  \item \textbf{Informed Consent:} A detailed discussion with the patient covers the indications, surgical procedure, expected outcomes, and potential risks and complications, including lymphedema, nerve injury, and seroma. The patient's understanding and consent are documented.
  \item \textbf{Pre-habilitation and Education:} Patients may receive education on postoperative arm exercises and lymphedema precautions. Smoking cessation counseling is offered to smokers.
\end{itemize}

\section*{Contraindications \& Precautions}
While axillary node clearance is a critical procedure for appropriate patients, certain conditions may contraindicate or necessitate precautions for its performance.

\textbf{Absolute Contraindications:}
\begin{itemize}
  \item \textbf{Uncontrolled active infection in the axilla:} Performing surgery through an infected field significantly increases the risk of wound infection and sepsis.
  \item \textbf{Severe, uncorrectable coagulopathy:} An inability to achieve adequate hemostasis poses an unacceptable risk of hemorrhage.
  \item \textbf{Patient refusal:} A competent patient's informed decision to decline surgery must be respected.
\end{itemize}

\textbf{Relative Contraindications \& Precautions:}
\begin{itemize}
  \item \textbf{Inoperable metastatic disease with poor prognosis:} In patients with widespread metastatic disease and a very limited life expectancy, the morbidity of surgery may outweigh any potential benefit. Palliative care or systemic therapy may be more appropriate.
  \item \textbf{Severe, uncompensated medical comorbidities:} Uncontrolled cardiac disease, severe pulmonary dysfunction, or other critical medical conditions that make general anesthesia or the stress of surgery excessively risky. Optimization of these conditions is paramount.
  \item \textbf{Extensive prior radiation to the axilla:} Previous radiation therapy can cause significant tissue fibrosis, making dissection more challenging and increasing the risk of complications such as lymphedema, wound healing issues, and nerve injury.
  \item \textbf{Pre-existing severe lymphedema of the ipsilateral arm:} Further axillary surgery can exacerbate existing lymphedema, leading to increased patient discomfort and functional impairment. Careful consideration and patient counseling are required.
  \item \textbf{High burden of nodal disease with expected poor response to surgery alone:} In such cases, neoadjuvant systemic therapy or primary radiation may be considered first to downstage the disease or as the definitive local treatment.
  \item \textbf{Anatomical variations or previous surgery:} Prior surgical interventions in the axilla can alter normal anatomy, making nerve and vessel identification more challenging and requiring increased surgical precision.
\end{itemize}
Precautions always include meticulous surgical technique, careful identification and preservation of nerves and vessels, and thorough hemostasis to minimize complications. A multidisciplinary tumor board discussion is often beneficial for complex cases.

\section*{Complications \& Risks}
Axillary node clearance, like any surgical procedure, carries a range of potential complications and risks, which can be broadly categorized as general surgical risks and procedure-specific risks.

\textbf{General Surgical Complications:}
\begin{itemize}
  \item \textbf{Bleeding and Hematoma:} Accumulation of blood in the surgical site, potentially requiring re-operation.
  \item \textbf{Infection:} Surgical site infection (SSI) of the wound or drain site, requiring antibiotics or drainage.
  \item \textbf{Pain:} Postoperative pain, which is usually manageable with analgesics but can occasionally be chronic.
  \item \textbf{Anesthesia Risks:} Adverse reactions to anesthetic agents, respiratory or cardiac complications, and nausea/vomiting.
  \item \textbf{Scarring:} Formation of a visible scar, which can sometimes be hypertrophic or keloid.
\end{itemize}

\textbf{Procedure-Specific Complications:}
\begin{itemize}
  \item \textbf{Lymphedema:} The most significant and feared long-term complication, characterized by chronic swelling of the arm, hand, and sometimes the trunk, due to disruption of lymphatic drainage. It can be debilitating and impact quality of life.
  \item \textbf{Seroma:} A collection of clear, serous fluid in the axillary dead space, which is very common and may require repeated aspiration or prolonged drain placement.
  \item \textbf{Nerve Injury:}
    \begin{itemize}
      \item \textbf{Intercostobrachial nerve injury:} Leads to numbness, tingling, or burning pain (dysesthesia) in the medial upper arm and axilla. This is common and often unavoidable.
      \item \textbf{Long thoracic nerve injury:} Results in weakness of the serratus anterior muscle, causing "winged scapula" (scapular protraction deficit) and difficulty raising the arm above the head.
      \item \textbf{Thoracodorsal nerve injury:} Causes weakness of the latissimus dorsi muscle, affecting arm adduction, extension, and internal rotation.
      \item \textbf{Brachial plexus injury:} Rare but severe, leading to significant motor and sensory deficits in the arm and hand.
    \end{itemize}
  \item \textbf{Shoulder Stiffness and Reduced Range of Motion:} Due to scar tissue formation, pain, or nerve injury, often requiring physiotherapy.
  \item \textbf{Axillary Web Syndrome (Cording):} Palpable and sometimes painful cords that develop in the axilla or down the arm, restricting movement.
  \item \textbf{Vascular Injury:} Rare but serious injury to the axillary artery or vein, potentially leading to significant hemorrhage or ischemia.
\end{itemize}

\section*{Postoperative Recovery Timeline}
The postoperative recovery timeline for axillary node clearance varies depending on whether it was performed alone or in conjunction with breast surgery, but generally follows a predictable course.

Immediately after surgery, patients are transferred to the Post-Anesthesia Care Unit (PACU) for close monitoring of vital signs, pain control, and assessment for immediate complications like bleeding. Most patients will stay in the hospital for one to two nights, especially if combined with a mastectomy. During the first 24-48 hours, pain is managed with oral or intravenous analgesics, and the surgical drain output is monitored. Patients are encouraged to begin gentle arm and shoulder exercises, as instructed by the surgical team or a physical therapist, to prevent stiffness and promote lymphatic flow. Early ambulation is also encouraged.

Over the first 1-2 weeks, the drain typically remains in place until output decreases to an acceptable level (e.g., <20-30 mL/day), at which point it is removed during an outpatient visit. Sutures or staples are usually removed around 10-14 days post-op. Patients are advised to avoid heavy lifting, strenuous activities, and repetitive arm movements on the affected side. Diet can usually be advanced as tolerated. Full range-of-motion exercises are gradually increased under guidance, and patients are educated on lymphedema precautions, including skin care, avoiding tight clothing, and protecting the arm from injury. Long-term recovery involves continued adherence to exercise programs and lymphedema surveillance. Most patients can return to light daily activities within 2-4 weeks, with full recovery and return to more strenuous activities potentially taking 6-12 weeks or longer, depending on individual healing and the extent of surgery.

\section*{Surgical Findings \& Pathology}
During axillary node clearance, the surgeon meticulously dissects the axillary contents, noting the presence and appearance of any enlarged or suspicious lymph nodes, the extent of their involvement, and their relationship to surrounding structures. Any gross evidence of extracapsular extension (cancer cells extending beyond the lymph node capsule) is also documented. The entire fibrofatty tissue containing the Level I and II lymph nodes is removed as a single specimen.

This resected specimen is then sent to the pathology laboratory. The pathologist's report is crucial for definitive cancer staging and treatment planning. It typically details:
\begin{itemize}
  \item \textbf{Total number of lymph nodes removed:} This indicates the completeness of the dissection.
  \item \textbf{Number of lymph nodes containing metastatic cancer:} This is a primary determinant of nodal stage.
  \item \textbf{Size of the largest metastatic deposit:} Provides prognostic information.
  \item \textbf{Presence or absence of extracapsular extension (ECE):} ECE indicates a higher tumor burden and is an adverse prognostic factor, often influencing recommendations for adjuvant radiation therapy.
  \item \textbf{Lymph node level involvement:} Confirmation of which levels (I, II) were involved.
  \item \textbf{Other pathological features:} Such as the presence of isolated tumor cells or micrometastases, though these are less common indications for full clearance.
\end{itemize}
These findings, combined with the primary tumor characteristics, determine the overall pathological stage of the breast cancer and guide subsequent adjuvant therapies.

\section*{Expected Postoperative Course}
A normal and successful postoperative course following axillary node clearance is characterized by a gradual resolution of pain, progressive restoration of arm and shoulder function, and the absence of significant complications. Initially, patients will experience incisional pain and discomfort in the axilla, which should be well-controlled with prescribed oral analgesics and progressively diminish over days to weeks. Swelling and bruising around the surgical site are common but should also subside.

The surgical drain, if placed, will typically be removed within the first week or two as fluid output decreases. Patients are expected to regain full or near-full range of motion in the affected arm and shoulder through adherence to a prescribed exercise regimen, usually within 4-8 weeks. Numbness or altered sensation in the medial arm due to intercostobrachial nerve manipulation is a common and often expected finding, which may improve over time but can be permanent. The wound should heal without signs of infection, and the scar will mature over several months. A successful course means the patient can resume normal daily activities, including work and hobbies, with minimal functional limitation and without developing chronic lymphedema or other debilitating complications. Regular follow-up with the surgical and oncology teams ensures that any deviations from this expected course are promptly addressed.

\section*{Postoperative Care \& Follow-up}
Comprehensive postoperative care and diligent follow-up are critical after axillary node clearance to monitor healing, manage potential complications, and coordinate further cancer treatment.

\begin{itemize}
  \item \textbf{Initial Follow-up Visit:} Typically scheduled 1-2 weeks post-surgery to assess wound healing, remove any remaining sutures or staples, and remove the surgical drain if output criteria are met.
  \item \textbf{Pathology Review:} A multidisciplinary team meeting, including the surgeon, oncologist, and radiation oncologist, reviews the final pathology report to finalize the cancer stage and formulate the adjuvant treatment plan (chemotherapy, radiation therapy, endocrine therapy).
  \item \textbf{Physiotherapy/Rehabilitation:} Referral to a physical therapist or occupational therapist is often recommended to guide specific exercises aimed at restoring full range of motion, preventing shoulder stiffness, and managing any early signs of axillary web syndrome or lymphedema.
  \item \textbf{Lymphedema Surveillance and Management:} Patients are educated on lifelong lymphedema precautions. Regular monitoring for arm swelling is crucial, with prompt referral to a specialized lymphedema therapist for early intervention (e.g., manual lymphatic drainage, compression garments) if symptoms develop.
  \item \textbf{Long-term Oncologic Follow-up:} Regular clinical examinations and appropriate imaging (e.g., mammography, breast MRI) continue according to established breast cancer surveillance guidelines to monitor for local, regional, or distant recurrence.
  \item \textbf{Warning Signs:} Patients are instructed to contact their surgical team immediately if they experience fever, increasing pain, spreading redness or warmth around the wound, excessive swelling of the arm, persistent drainage from the wound, or any signs of wound breakdown.
\end{itemize}
Ongoing psychological support and attention to quality of life issues are also important components of long-term care.

\section*{Epidemiology}
Breast cancer is the most common cancer among women globally, and the axillary lymph nodes are the most frequent site of regional metastatic spread. Historically, axillary node clearance was a standard component of breast cancer surgery, leading to a high incidence of associated morbidity. However, the epidemiology of ALND has dramatically shifted with the widespread adoption of sentinel lymph node biopsy (SLNB) since the late 1990s.

Currently, the majority of early-stage breast cancer patients are found to be node-negative by SLNB, thus avoiding full axillary clearance. Consequently, the incidence of ALND has significantly decreased, now being performed in a more select group of patients, primarily those with clinically positive nodes, multiple positive sentinel nodes, or residual disease after neoadjuvant therapy. The risk of lymphedema after ALND is substantial, affecting 10-40\% of patients, depending on the extent of dissection, adjuvant radiation, and patient factors. This morbidity has been a major driver in the de-escalation of axillary surgery. While ALND provides crucial staging information and local control, its impact on overall breast cancer mortality has been increasingly questioned in the era of effective systemic therapies, leading to its more targeted application.

\section*{Outcomes \& Prognosis}
The outcomes of axillary node clearance are primarily assessed in terms of local-regional disease control, accurate pathological staging, and patient morbidity. ALND provides excellent local control of the axilla, effectively removing macroscopic and microscopic nodal disease, thereby reducing the risk of regional recurrence. The detailed pathological information obtained (number of positive nodes, extracapsular extension) is critical for accurate staging, which is a powerful prognostic factor for overall survival and guides the intensity of adjuvant systemic and radiation therapies.

However, the impact of ALND on overall survival has been refined by landmark clinical trials. The ACOSOG Z0011 trial demonstrated that in patients with T1/T2 breast cancer, 1-2 positive sentinel nodes, and planned whole-breast radiation, omitting ALND did not compromise overall or disease-free survival compared to performing a complete axillary dissection. This trial significantly altered practice, reducing the number of ALNDs performed. For patients who do undergo ALND, the prognosis is strongly influenced by the underlying biology of the breast cancer, the number of involved nodes, and the effectiveness of adjuvant therapies. Patients with fewer positive nodes and favorable tumor biology generally have a better prognosis. The main long-term morbidity, lymphedema, significantly impacts patient-reported outcomes and quality of life, underscoring the importance of judicious patient selection for this procedure.

\section*{References}
\begin{enumerate}
  \item Giuliano AE, Hunt KK, Morton DL, et al. Axillary Dissection vs. No Axillary Dissection in Patients with T1-T2 N0-N1 M0 Breast Cancer and 1-2 Positive Sentinel Nodes: A Randomized Clinical Trial. JAMA. 2011;305(6):569-575.
  \item Schwartz SI, Brunicardi FC, Andersen DK, et al. Schwartz's Principles of Surgery. 11th ed. McGraw-Hill Education; 2019.
  \item National Comprehensive Cancer Network (NCCN). NCCN Clinical Practice Guidelines in Oncology: Breast Cancer. Version 1.2024.
  \item MedlinePlus. Axillary lymph node dissection. U.S. National Library of Medicine; 2024.
\end{enumerate}

\clearpage